\section*{Chapter \ref{ch:b2:dipole-radiation} --- Dipole Radiation and Scattering}

\begin{solution}{exo:b2:dipole-radiation:1}
$p_0 = I_0\ell/\omega = \qty{3.2e-9}{C.m}$; $E = \mu_0\omega^2p_0\sin\theta/4\pi r$: \qty{0.13}{V/m}
at $\theta = 90^\circ$, \qty{0.063}{V/m} at \qty{30}{\degree}; $\mathcal P = \mu_0\omega^4p_0^2/12\pi c =
\qty{175}{W}$; $R_{\text{rad}} = 2\mathcal P/I_0^2 = \qty{88}{\ohm}$.
\end{solution}

\begin{solution}{exo:b2:dipole-radiation:2}
$790(\ell/\lambda)^2$: \qty{8.8e-5}{\ohm}, \qty{8.8e-3}{\ohm}, \qty{0.88}{\ohm}, \qty{88}{\ohm}. Copper:
$0.026$, $0.084$, $0.27$, \qty{0.84}{\ohm} (\omterm{thm:b2:waves-in-media:skin}{skin effect}): efficiencies $0.3\%$, $10\%$,
$77\%$, $99\%$. At long wavelengths only a mast that is a fair fraction of
$\lambda$ has a \omterm{ex:b2:dipole-radiation:antenna}{radiation resistance} above its losses.
\end{solution}

\begin{solution}{exo:b2:dipole-radiation:3}
$400 : 550 : 700 = 9.4 : 2.6 : 1$. $(10/3)^4 = 120$: the \qty{3}{cm} radar
sees small drops a hundred times better; the \qty{10}{cm} one looks
through the rain to what is behind.
\end{solution}

\begin{solution}{exo:b2:dipole-radiation:4}
$\sin^2\chi/(1 + \cos^2\chi)$: $0.14$, $0.60$, $1$, $0.60$. Point \qty{90}{\degree} from the
Sun (a great circle across the sky); with the Sun overhead, the whole
horizon is at \qty{90}{\degree} and polarized along it. Multiple scattering,
the ground's reflection and the molecules' anisotropy add unpolarized
light: about $75\%$ at best.
\end{solution}

\begin{solution}{exo:b2:dipole-radiation:5}
(a) $\int\sin^2\theta\,\dd\Omega = 8\pi/3$. (b) $\int_{60^\circ}^{120^\circ}\sin^3\theta\,\dd\theta/
(4/3) = 0.917/1.333 = 69\%$. (c) $\sin^2\theta = \tfrac12$: $\theta = \qty{45}{\degree}$, a
beam \qty{90}{\degree} wide. (d) Maximum $1$ against the average $2/3$: $1.5$.
\end{solution}

\begin{solution}{exo:b2:dipole-radiation:6}
(a) $p/4\pi\varepsilon_0r^3 = \omega^2p/4\pi\varepsilon_0c^2r$ at $r_0 = c/\omega = \lambda/2\pi$. (b) \qty{48}{m}:
at \qty{10}{m} the quasi-static field, at \qty{10}{km} the radiated one. (c)
$r_0 = \qty{4.8}{cm}$: the head is in the near zone. (d) Near-zone $\vect E$
and $\vect B$ are in quadrature: the \omterm{thm:b2:poynting-vector:poynting}{Poynting vector} averages to zero,
energy is stored and returned, as in a capacitor.
\end{solution}

\begin{solution}{exo:b2:dipole-radiation:7}
(a) $8\pi r_e^2/3 = \qty{6.65e-29}{m^2}$. (b) $n_e\sigma_TL = 10^{14} \times 6.65 \times 10^{-29}
\times 10^9 = 7 \times 10^{-6}$: a millionth of the photosphere, colourless because
\omterm{prop:b2:dipole-radiation:rayleigh}{Thomson scattering} is, and drowned by the sky's scattered sunlight
except at eclipse. (c) \qty{10}{keV} $\gg$ \qty{300}{eV}: the electrons respond as
free, $6\sigma_T = \qty{4e-28}{m^2}$. (d) Photoelectric absorption grows as $Z^4$
and singles out calcium; scattering is the same per electron in bone
and flesh.
\end{solution}

\begin{solution}{exo:b2:dipole-radiation:8}
(a) $\tau = 0.36$, $0.10$, $0.038$: $T = 0.70$, $0.90$, $0.96$. (b) $\tau \times 38 = 13.6$,
$3.8$, $1.45$: $T = 1 \times 10^{-6}$, $0.02$, $0.23$; red over blue $\sim 2 \times 10^5$. (c)
The Earth's atmosphere bends sunset light into the shadow, reddened
by the same scattering. (d) Near the horizon the path is long: the
blue is scattered out of the scattered light itself, and multiple
scattering adds white.
\end{solution}

\begin{solution}{exo:b2:dipole-radiation:9}
(a) Along the line of the pair the path difference $\lambda/2$ gives
cancellation; broadside, addition. (b) Antiphase: the reverse ---
radiation along the line (end-fire), none broadside. (c) Toward $+x$:
path delay $-\pi/2$ plus feed phase $+\pi/2$, in phase; toward $-x$:
$-\pi$, cancelling: a cardioid, one-sided. (d) Masts shape the coverage
and protect neighbours; a phased array steers by changing the feed
phases, in microseconds.
\end{solution}

\begin{solution}{exo:b2:dipole-radiation:10}
(a) $\langle\mathcal P\rangle = e^2\omega_0^4x_0^2/12\pi\varepsilon_0c^3$, $W = \tfrac12m\omega_0^2x_0^2$:
$\Gamma = \mathcal P/W = e^2\omega_0^2/6\pi\varepsilon_0mc^3$. (b) $\omega_0 = \qty{3.2e15}{rad/s}$: $\Gamma =
\qty{6.4e7}{s^{-1}}$ --- the value used. (c) \qty{16}{ns}; $\Delta\lambda = \lambda^2\Gamma/2\pi c =
\qty{1.2e-14}{m}$, a hundredth of a picometre (\qty{10}{MHz}). (d) $Q = 5 \times 10^7$.
\end{solution}

\begin{solution}{exo:b2:dipole-radiation:11}
(a) \qty{75}{m}; $I_0 = \sqrt{2\mathcal P/R} = \qty{53}{A}$. (b) $\langle\Pi\rangle = 1.5 \times
5 \times 10^4/2\pi \times 10^8 = \qty{1.2e-4}{W/m^2}$, $E_0 = \sqrt{2\Pi/\varepsilon_0c} = \qty{0.30}{V/m}$.
(c) \qty{0.30}{V}. (d) The return currents flow in the ground; a resistive
soil adds losses and spoils the mirror; buried radials give them a
copper path.
\end{solution}

\begin{solution}{exo:b2:dipole-radiation:12}
(a) $p_0 = \alpha E_0$ gives $\sigma = \alpha^2\omega^4/6\pi\varepsilon_0^2c^4$; with $\alpha = \varepsilon_0(n^2
- 1)/N$ and $\omega = 2\pi c/\lambda$: $\sigma = 8\pi^3(n^2 - 1)^2/3N^2\lambda^4$. (b) $n^2 - 1 =
5.8 \times 10^{-4}$: $\sigma = \qty{4.9e-31}{m^2}$, mean free path $1/N\sigma = \qty{80}{km}$. (c)
$8/80 = 0.1$. (d) $\sigma$ is measured from the sky's \omterm{def:b2:dispersion-wave-packets:complex}{attenuation}, $n$ is
known: $N$ follows, hence Avogadro's number.
\end{solution}

\begin{solution}{pb:b2:dipole-radiation:1}
\textbf{1.} $\lambda = \qty{300}{m}$, $h = \qty{75}{m}$; $I_0 = \sqrt{2 \times 10^5/36} = \qty{75}{A}$;
$V = RI_0 = \qty{2.7}{kV}$.

\textbf{2.} $\ell_{\text{eff}} = \qty{48}{m}$: $p_0 = 75 \times 48/6.3 \times 10^6 = \qty{5.7e-4}{C.m}$.

\textbf{3.} $\langle\Pi\rangle = 1.5 \times 10^5/2\pi r^2$: \qty{2.4e-4}{W/m^2} and \qty{2.4e-6}{W/m^2};
$E_0 = \qty{0.42}{V/m}$ and \qty{0.042}{V/m}.

\textbf{4.} Whip: \qty{42}{mV}. Rod: $N\mu_{\text{eff}}A\omega B_0 = 10^4 \times 10^{-4} \times 6.3 \times
10^6 \times 1.4 \times 10^{-10} = \qty{0.9}{mV}$ --- less, but compact and selective.

\textbf{5.} $A = 1.5\lambda^2/4\pi = \qty{1.1e4}{m^2}$: $P = \Pi A = \qty{26}{mW}$.

\textbf{6.} $36/38 = 95\%$; \qty{5}{kW} of heat.

\textbf{7.} $r/\lambda = 33$ at \qty{10}{km}: \omterm{thm:b2:dipole-radiation:field}{radiation zone}; at \qty{100}{m},
$r \approx \lambda/3$: the transition region, where the quasi-static field still
counts.

\textbf{8.} By day the D layer absorbs the wave that would otherwise
reach the reflecting layers; at night it is gone and the sky wave
carries the station hundreds of kilometres.

\textbf{9.} $\sigma(450) = \qty{1.0e-30}{m^2}$, $\sigma(650) = \qty{2.4e-31}{m^2}$: ratio $4.4$.

\textbf{10.} Mean free paths $39$, $87$, \qty{170}{km}; $\tau = 0.21$, $0.09$, $0.05$:
$19\%$, $9\%$, $5\%$ scattered.

\textbf{11.} $N\sigma I = 2.5 \times 10^{25} \times 4.6 \times 10^{-31} \times 300 = \qty{3.5}{mW/m^3}$,
into $4\pi$ with the pattern $(1 + \cos^2\chi)$ for \omterm{def:b2:plane-waves-polarization:states}{natural light}.

\textbf{12.} The column scatters $3.5 \times 10^{-3} \times 8000 = \qty{28}{W}$ of its
\qty{300}{W} ($9\%$); spread over the hemisphere, the sky delivers some
\qty{4}{W.m^{-2}.sr^{-1}}, against the Sun's $300/6.8 \times 10^{-5} = \qty{4e6}{W.m^{-2}.sr^{-1}}$:
a millionth per steradian --- the right order.

\textbf{13.} Fully polarized at \qty{90}{\degree} in single scattering; multiple
scattering and the ground reduce it to some $75\%$; bees and ants read
the pattern to find the Sun behind clouds.

\textbf{14.} $\tau \times 38$: blue $8$ ($T = 3 \times 10^{-4}$), red $1.8$ ($T = 0.17$): a deep
red Sun; overhead, the light arrives after a short path and is blue.

\textbf{15.} No atmosphere, nothing to scatter: the sky is black beside
the Sun.

\textbf{16.} Micrometre dust scatters like Mie particles: forward and
red, giving a butterscotch sky; near the setting Sun the
forward-scattered light is blue.

\textbf{17.} $8\pi r_e^2/3 = \qty{6.65e-29}{m^2}$.

\textbf{18.} $a = v^2/2d = \qty{6.4e21}{m/s^2}$; $\mathcal P = e^2a^2/6\pi\varepsilon_0c^3 =
\qty{2e-10}{W}$ during $2d/v = \qty{2.5e-14}{s}$: \qty{6e-24}{J}, $10^{-10}$ of
\qty{100}{keV} --- the real yield is near $1\%$, because the braking happens
in violent close encounters with nuclei, not in a gentle uniform
deceleration.

\textbf{19.} $hc/E = \qty{12.4}{pm}$.

\textbf{20.} \qty{2}{kW} in the beam, about \qty{20}{W} of X-rays; two kilowatts
of heat in a spot --- the anode rotates to spread it.

\textbf{21.} $a = v^2/r = \qty{9e22}{m/s^2}$, $\mathcal P = \qty{5e-8}{W}$; \qty{13.6}{eV} gone
in \qty{5e-11}{s}.

\textbf{22.} Classical charges in bound motion radiate continuously;
quantum mechanics has stationary states that do not, and a lowest one
with nowhere to fall.

\textbf{23.} $a = c^2/R = \qty{9e14}{m/s^2}$: \qty{5e-24}{W} by Larmor, $\gamma^4 \sim 10^{16}$
times more in reality: tens of kilowatts of synchrotron light from a
stored beam.

\textbf{24.} $\sigma_TI = \qty{6.7e-26}{W}$; $1.5 \times 10^{25}$ electrons for a watt.

\textbf{25.} Mast: \qty{5.7e-4}{C.m} at \qty{1}{MHz}, \qty{100}{kW}, a vertical
doughnut. Molecule: $\sim$\qty{1e-35}{C.m} at \qty{500}{THz}, \qty{1e-28}{W},
a doughnut about the field. Free electron: Thomson, \qty{1e-25}{W} in
sunlight, the same pattern.
\end{solution}
